\section{The functional equation makes the problem singular}\label{es:sec:singular}

\subsection{Pairs, not zeros}

Theorem~\ref{es:thm:concentration} treats each zero as a free parameter $(\gamma,a)$. The zeros are not free: $\xi(s)=\xi(1-s)$, so a zero $\rho=\tfrac12+\delta+i\gamma$ with $\delta>0$ comes with its partner $1-\bar\rho=\tfrac12-\delta+i\gamma$ at the same ordinate. At level $\theta$ the pair has widths $a_0\mp\delta$ with $a_0=\theta-\tfrac12$, and its two poles $e^{\Delta(-(a_0\mp\delta)+i(\gamma-c))}$ merge as $\delta\to0$ into one pole of doubled amplitude, the contribution of a double zero on the line. The pair enters the samples with one amplitude common to its two members, so that the contribution of a pair is $A\,[f(a_0-\delta)+f(a_0+\delta)]$ with one free complex amplitude $A$, and $\delta=0$ is a zero on the line of free amplitude $2A$. The right alternative to ``this zero is on the line'' is therefore not ``this zero has width $a_0+\varepsilon$'', which the functional equation forbids, but ``this zero is a pair at $\tfrac12\pm\delta$''. The two alternatives behave very differently.

\begin{lemma}[evenness]\label{es:lem:even}
In the notation of Theorem~\ref{es:thm:taper} at level $\theta$, with $0\le\delta<a_0$, write the contribution of one zero of width $a$ at ordinate $\gamma$ to the samples as $f(a)_j=\exp(-q^2/2\ell^2+i\omega_jq)$ with $q=(\gamma-c)+ia$. The contribution of the pair at $\tfrac12\pm\delta$ is $f(a_0-\delta)+f(a_0+\delta)$, an even entire function of $\delta$, hence an entire function $g(\eta)$ of $\eta=\delta^2$, with
\[
\begin{gathered}g(\eta)=2f(a_0)+\eta\,t+O(\eta^2),\\ t_j=f''(a_0)_j=f(a_0)_j\Bigl[\Bigl(\omega_j+\frac{iq_0}{\ell^2}\Bigr)^2+\frac1{\ell^2}\Bigr],\quad q_0=(\gamma-c)+ia_0 .\end{gathered}
\]
\end{lemma}

\begin{proof}
Swapping the two members of the pair changes $\delta$ to $-\delta$ and leaves the sum unchanged, so the sum is even; it is entire in $\delta$ because $f$ is entire in $a$. An even entire function is an entire function of $\delta^2$. The coefficient of $\delta^2$ is $f''(a_0)$, and $\partial_a\log f=-(\omega_j+iq/\ell^2)$, $\partial_a^2\log f=1/\ell^2$, which gives $t$.
\end{proof}

With the amplitudes of Theorem~\ref{es:thm:taper}, the limit $\delta\to0$ of a pair is a double zero, $2f(a_0)$, and a simple zero, $f(a_0)$, is not nested in the pair model. The singular analysis therefore compares a double zero with a pair, or, in the common amplitude model, a zero of free amplitude with a pair of one free common amplitude; for a simple zero the amplitude is $A^0=\tfrac12$ and the tangent is $t=f''(a_0)/2$. Freeing the amplitudes of the two members independently discards the multiplicity: the contribution of the pair is then not even in $\delta$, and the joint Jacobian is rank deficient at $\delta=0$. We use the common amplitude model throughout this section, and the statements about ``a zero'' below are statements in the free amplitude model.

So the data depend on the displacement $\delta$ only through $\delta^2$: the Fisher information for $\delta$ vanishes at $\delta=0$, and the model is singular there, in the way a mixture with two merging components is singular \cite{Chen95,Nguyen13}. In the parameter $\eta\ge0$ it is regular, one sided at $\eta=0$, when $t_\perp\ne0$ (Theorem~\ref{es:thm:singular}).

\begin{theorem}[the rate for the real part]\label{es:thm:singular}
Fix a window, a level $\theta>\tfrac12$ with $a_0=\theta-\tfrac12$, the samples $\omega_j=\omega_0+j\Delta$, $0\le j<J$, and $f$ as in Lemma~\ref{es:lem:even}. Let one ordinate $\gamma$ carry either a zero on the line or a pair at $\tfrac12\pm\delta$, $0\le\delta\le\delta_+<a_0$, entering the samples as $A\,g(\eta;\gamma)$, where $g(\eta;\gamma)=f(a_0-\sqrt\eta)+f(a_0+\sqrt\eta)$, $\eta=\delta^2$, and $A$ is one complex amplitude common to the two members; let the other $M$ zeros of the window enter as $\sum_{m\le M}A_mz_m^{\,j}$ as in Theorem~\ref{es:thm:concentration}, labelled by increasing ordinate. Write $\vartheta'=(\gamma,A,(\gamma_m,a_m,A_m)_{m\le M})$, of real dimension $D'=3+4M$, and $\Phi(\eta,\vartheta')\in\C^J=\R^{2J}$ for the noiseless samples, and assume:
\begin{enumerate}
\item[(H1)] $Y=\Phi(\eta^0,\vartheta'^0)+b+\tau\xi$, where $\xi\in\C^J$ has independent standard normal real and imaginary parts, as in Theorem~\ref{es:thm:concentration}, so that $\tau$ is the standard deviation of each real coordinate, and $b\in\C^J$ is deterministic with $\|b\|_2\le B\tau$;
\item[(H2)] the parameter set is $\Theta=[0,\eta_+]\times K'$ with $\eta_+=\delta_+^2$ and $K'$ compact, $\eta^0\in[0,\eta_+)$, $\vartheta'^0$ is interior to $K'$, and $\Phi$ is injective on $\Theta$;
\item[(H3)] the real $2J\times D'$ Jacobian $T'$ of $\vartheta'\mapsto\Phi(\eta^0,\vartheta')$ at $\vartheta'^0$ has rank $D'$; with $t=\partial_\eta\Phi(\eta^0,\vartheta'^0)$, which is $A^0f''(a_0)$ when $\eta^0=0$, $P'$ the orthogonal projection of $\R^{2J}$ onto the range of $T'$ and $t_\perp=(I-P')t$, we have $\|t_\perp\|>0$;
\item[(H4)] the prior puts mass $1-p\in[0,1)$ at $\eta=0$ and has density $p\,\pi_\eta(\eta)$ on $(0,\eta_+]$, $\pi_\eta(\eta)=\pi(\sqrt\eta)/2\sqrt\eta$ with $\pi$ continuous and positive on $[0,\delta_+]$, and, independently, $\vartheta'$ has a density $p'$ bounded on $K'$, continuous and positive at $\vartheta'^0$.
\end{enumerate}
Let $r_\tau=\tau/\|t_\perp\|$ and let $\tau\to0$ with $B$ and $\vartheta'^0$ fixed, and with $\eta^0$ fixed in (i) and (ii). Then:
\begin{enumerate}
\item[(i)] for every $M_\tau\to\infty$, $\Pi(|\eta-\eta^0|>M_\tau r_\tau\mid Y)\to0$ in probability; hence $\Pi(\delta>(M_\tau r_\tau)^{1/2}\mid Y)\to0$ if $\delta^0=0$, and $\Pi(|\delta-\delta^0|>M_\tau r_\tau/\delta^0\mid Y)\to0$ if $\delta^0>0$;
\item[(ii)] let $\eta^0=0$ and, for $0<\delta_1\le\delta_+$, let $D(\delta_1)=\min_{\vartheta'\in K'}\|\Phi(\delta_1^2,\vartheta'^0)-\Phi(0,\vartheta')\|$, attained at $\vartheta'^*$; then $D(\delta_1)=\delta_1^2\|t_\perp\|+O(\delta_1^4)$, the Kullback Leibler divergence of the law of $Y$ under $(\delta_1^2,\vartheta'^0)$ from the law under $(0,\vartheta'^*)$ is $D(\delta_1)^2/2\tau^2$, and for $\delta_1=c\,r_\tau^{1/2}$ the two error probabilities of any test of $\delta=0$ against $\delta=\delta_1$, whatever the method, sum to at least $1-c^2/2-O(\tau)$;
\item[(iii)] let $\eta^0=v\,r_\tau$ with $v\ge0$ fixed ($v=0$ is the null), let $\|b\|_2=o(\tau)$, take $t_\perp$ at $\eta=0$, and let $Z=\langle t_\perp,r\rangle/(\tau\|t_\perp\|)$ with $r$ the residual of the least squares fit of $\Phi(0,\cdot)$ over $K'$; then $Z=v+\zeta_\perp+o_P(1)$ with $\zeta_\perp$ standard Gaussian, so the one sided matched filter $Z>z_\alpha$ has asymptotic level $\alpha$ and asymptotic power $1-N(z_\alpha-v)$, $N$ the standard normal distribution function, which is $\tfrac12$ at $\delta^*=(z_\alpha r_\tau)^{1/2}$; and twice the logarithm of the generalised likelihood ratio of $\eta\ge0$ against $\eta=0$ is $\max(Z,0)^2+o_P(1)$, so that for $\alpha<\tfrac12$ the generalised likelihood ratio test at level $\alpha$ is asymptotically the matched filter.
\end{enumerate}
\end{theorem}

\begin{remark}\label{es:rem:singular}
(H2) and (H3) hold, and $\|t_\perp\|>0$ is automatic when $\eta^0=0$, if $J\ge2M+4$, $2R\Delta<2\pi$ as in Theorem~\ref{es:thm:concentration}, the amplitudes in $K'$, $A$ included, are bounded away from $0$, the ordinate $\gamma$ ranges over an interval $I$ of length less than $2\pi/\Delta$, and the poles $z_m$ keep a distance at least $s>0$ from one another and from the set $S=\{e^{\Delta(-a+i(\gamma'-c))}:|a-a_0|\le\delta_+,\ \gamma'\in I\}$ that contains the poles of the pair. Indeed $f(a)_j=f(a)_0\,z(a)^j$ with $z(a)=e^{\Delta(-a+i(\gamma-c))}$ and $f(a)_0\ne0$, so two parameters with the same samples give two exponential sums whose difference has at most $2M+4\le J$ distinct poles and vanishes for $0\le j<J$; the Vandermonde matrix of these poles has full column rank, so the two sums have the same poles with the same coefficients. The poles in $S$ are those of the pair: one pole means $\eta=0$, two poles of moduli $e^{-\Delta(a_0\mp\sqrt\eta)}$ give $\eta>0$, their common argument gives $\gamma$, and the coefficients divided by $f(a_0\mp\sqrt\eta)_0$ give $A$; the remaining poles give the $(\gamma_m,a_m,A_m)$, the ordinates by $2R\Delta<2\pi$ and the labels by their order. For (H3) at $\eta^0=0$, the real columns of $T'$ lie in the complex span of $(z_m^{\,j})_j$, $(jz_m^{\,j})_j$, $m\le M$, and $(z_0^{\,j})_j$, $(jz_0^{\,j})_j$, $z_0=z(a_0)$: four real columns per zero, $z_m^{\,j}$, $iz_m^{\,j}$, $iA_m\Delta jz_m^{\,j}$, $-A_m\Delta jz_m^{\,j}$, and three for the pair, independent over $\R$ because the confluent Vandermonde matrix of these $2M+2$ complex columns has full rank for $J\ge2M+2$. Finally $t=A^0f''(a_0)$ and $\partial_a^2z(a)^j=\Delta^2j^2z(a)^j$, so $t$ has the component $A^0\Delta^2f(a_0)_0\,(j^2z_0^{\,j})_j$ with a nonzero coefficient, and the confluent Vandermonde matrix with the further column $(j^2z_0^{\,j})_j$ has full rank for $J\ge2M+3$; so $t$ is not in the complex span of the columns of $T'$, a fortiori not in their real span, and $t_\perp\ne0$. With a prior on $M$ as in Theorem~\ref{es:thm:concentration} (amplitudes bounded below, $M\le M_{\max}$) and $J\ge2M_{\max}+4$, Step~1 of the proof extends across $M$ by the Vandermonde argument of that theorem, and (i) holds unchanged. Like Theorem~\ref{es:thm:concentration}, the theorem needs the bias below the noise and no aliasing, so by Remark~\ref{es:rem:settings} it does not cover the runs on the zeros of $\zeta$ in this chapter; it covers data simulated from the model with $b=0$ in a window of radius below $\pi/\Delta$.
\end{remark}

\begin{proof}
Write $\vartheta=(\eta,\vartheta')$, $\vartheta^0=(\eta^0,\vartheta'^0)$, $u(\vartheta)=\Phi(\vartheta)-\Phi(\vartheta^0)$ and $n=b+\tau\xi$, so that the log likelihood ratio of $\vartheta$ to $\vartheta^0$ is $\ell(\vartheta)=(2\langle u,n\rangle-\|u\|^2)/2\tau^2$, with $\langle\cdot,\cdot\rangle$ the real inner product of $\R^{2J}$. By Lemma~\ref{es:lem:even}, $\Phi$ is real analytic on a neighbourhood of $\Theta$, including $\eta<0$, so there are $L$ and $C_2$ with $\|u(\vartheta)\|\le L|\vartheta-\vartheta^0|$, $\|J_0\|\le L$ and $\|u(\vartheta)-J_0(\vartheta-\vartheta^0)\|\le C_2|\vartheta-\vartheta^0|^2$ on $\Theta$, where $J_0=[t,T']$.

\emph{Step 1 (separation).} By (H3), $t$ is not in the range of $T'$, so $J_0$ has full column rank and $\lambda=\sigma_{\min}(J_0)>0$; hence $\|u(\vartheta)\|\ge\lambda|\vartheta-\vartheta^0|/2$ for $|\vartheta-\vartheta^0|\le r_1=\lambda/2C_2$. By (H2), continuity and compactness, $g_1=\min\{\|u(\vartheta)\|:\vartheta\in\Theta,\ |\vartheta-\vartheta^0|\ge r_1\}>0$. With $\lambda_1=\min(\lambda/2,g_1/r_1)$, $\|u(\vartheta)\|\ge\lambda_1\min(|\vartheta-\vartheta^0|,r_1)$ on $\Theta$; note that $\lambda_1\le\lambda/2\le\sigma_{\min}(T')/2$.

\emph{Step 2 (localisation).} Let $E_\tau=\{\|\xi\|\le\sqrt{2J}+2(\log1/\tau)^{1/2}\}$, of probability at least $1-\tau^2$ by Gaussian concentration. On $E_\tau$, $\|n\|\le\tau\Lambda_\tau$ with $\Lambda_\tau=B+\sqrt{2J}+2(\log1/\tau)^{1/2}$, hence $\ell\le(2\|u\|\tau\Lambda_\tau-\|u\|^2)/2\tau^2$ everywhere and $\ell\ge-L\Lambda_\tau-L^2/2$ on $B_\tau=\{\vartheta\in\Theta:|\vartheta-\vartheta^0|\le\tau\}$. By (H4) the prior has a density bounded below by a positive constant on $\Theta$ near $\vartheta^0$, and $B_\tau$ contains at least half a ball, so $\Pi(B_\tau)\ge c_0\tau^{1+D'}$ and $\int e^\ell\,d\Pi\ge c_0\tau^{1+D'}e^{-L\Lambda_\tau-L^2/2}$. Let $\varepsilon_\tau=K\tau(\log1/\tau)^{1/2}$ with $\lambda_1K>\max(8,2(1+D')^{1/2})$. For small $\tau$, $\varepsilon_\tau\le r_1$ and $\lambda_1\varepsilon_\tau\ge4\tau\Lambda_\tau$; as $2x\tau\Lambda_\tau-x^2\le-x^2/2$ for $x\ge4\tau\Lambda_\tau$, $\ell\le-\lambda_1^2\varepsilon_\tau^2/4\tau^2$ on $\{|\vartheta-\vartheta^0|>\varepsilon_\tau\}$. Hence, on $E_\tau$,
\[
\Pi\bigl(|\vartheta-\vartheta^0|>\varepsilon_\tau\bigm|Y\bigr)\le c_0^{-1}\exp\Bigl(-\Bigl(\frac{\lambda_1^2K^2}4-1-D'\Bigr)\log\frac1\tau+L\Lambda_\tau+\frac{L^2}2\Bigr)\to0 .
\]
The radius $\varepsilon_\tau$ only localises the posterior; the rate comes from Steps 3 and 4.

\emph{Step 3 (linearisation).} On $U_\tau=\{\vartheta\in\Theta:|\vartheta-\vartheta^0|\le\varepsilon_\tau\}$ write $\vartheta-\vartheta^0=(h_\eta,h')$ and $k=h'+h_\eta w$ with $w=(T'^{\mathsf T}T')^{-1}T'^{\mathsf T}t$, so that $J_0(\vartheta-\vartheta^0)=h_\eta t_\perp+T'k$, the two terms being orthogonal. Then $\ell=\ell_1(h_\eta)+\ell_2(k)+\rho$ with
\[
\ell_1(x)=\frac{2x\langle t_\perp,n\rangle-x^2\|t_\perp\|^2}{2\tau^2},\qquad\ell_2(k)=\frac{2\langle T'k,n\rangle-\|T'k\|^2}{2\tau^2},
\]
and, on $E_\tau$, $|\rho|\le(2C_2\varepsilon_\tau^2\tau\Lambda_\tau+2LC_2\varepsilon_\tau^3+C_2^2\varepsilon_\tau^4)/2\tau^2=O(\tau(\log1/\tau)^{3/2})$ uniformly on $U_\tau$.

\emph{Step 4 (the marginal of $\eta$).} Let $W=\langle t_\perp,n\rangle/(\tau\|t_\perp\|)=\zeta_\perp+\langle t_\perp,b\rangle/(\tau\|t_\perp\|)$, where $\zeta_\perp=\langle t_\perp,\xi\rangle/\|t_\perp\|$ is standard Gaussian and the second term is at most $B$ in modulus; so $W=O_P(1)$ and $\ell_1(vr_\tau)=Wv-v^2/2$. Let $I_2=\int_{\R^{D'}}e^{\ell_2(k)}\,dk$, a Gaussian integral centred at $\hat k=(T'^{\mathsf T}T')^{-1}T'^{\mathsf T}n$ with covariance $\tau^2(T'^{\mathsf T}T')^{-1}$; on $E_\tau$, $|\hat k|\le\tau\Lambda_\tau/\sigma_{\min}(T')\le\varepsilon_\tau/4$ for small $\tau$, because $\sigma_{\min}(T')K\ge2\lambda_1K>16$ and $\Lambda_\tau\le4(\log1/\tau)^{1/2}$ for small $\tau$. Put $S_\tau=U_\tau\cap\{|h_\eta|>M_\tau r_\tau\}$ and $S'_\tau=U_\tau\cap\{|h_\eta|\le r_\tau\}$. The change of variables $(h_\eta,h')\mapsto(h_\eta,k)$ has Jacobian $1$. Bounding each $k$ integral by $I_2$,
\[
\int_{S_\tau}e^\ell\,d\Pi\le e^{o(1)}\sup_{K'}p'\;I_2\int_{|x|>M_\tau r_\tau}e^{\ell_1(x)}\,p\,\pi_\eta(\eta^0+x)\,dx .
\]
For $|h_\eta|\le r_\tau$ the $k$ section of $U_\tau$ contains the ball of radius $\varepsilon_\tau/2$ about $\hat k$, since there $|h'|\le|\hat k|+\varepsilon_\tau/2+r_\tau|w|\le(\tfrac34+o(1))\varepsilon_\tau$, and the Gaussian gives this ball the mass $(1-o(1))I_2$, its radius being large against $\tau/\sigma_{\min}(T')$; so
\[
\int_{S'_\tau}e^\ell\,d\Pi\ge e^{-o(1)}(1-o(1))\inf_{U_\tau}p'\;I_2\int_{|x|\le r_\tau}e^{\ell_1(x)}\,p\,\pi_\eta(\eta^0+x)\,dx .
\]
If $\eta^0>0$, then $\pi_\eta(\eta^0+x)=\pi_\eta(\eta^0)(1+o(1))$ for $|x|\le\varepsilon_\tau$, and with $x=vr_\tau$ the ratio of the right sides is at most $(1+o(1))\int_{|v|>M_\tau}e^{Wv-v^2/2}dv\big/\int_{|v|\le1}e^{Wv-v^2/2}dv$, which tends to $0$ in probability because $W=O_P(1)$; the atom at $\eta=0$ lies outside $U_\tau$ and is covered by Step 2. If $\eta^0=0$, then $x=\eta\ge0$, $\pi_\eta(\eta)=\pi(0)(1+o(1))/2\sqrt\eta$ on $(0,\varepsilon_\tau]$, the factor $\pi(0)r_\tau^{1/2}/2$ cancels, and the ratio is at most $(1+o(1))\int_{M_\tau}^\infty v^{-1/2}e^{Wv-v^2/2}dv\big/\int_0^1v^{-1/2}e^{Wv-v^2/2}dv\to0$; the atom at $\eta=0$ lies in $S'_\tau$ and only increases the denominator. As $\Pi(|\eta-\eta^0|>M_\tau r_\tau\mid Y)\le\Pi(U_\tau^c\mid Y)+\Pi(S_\tau\mid Y)/\Pi(S'_\tau\mid Y)$, this proves (i) for $\eta$; the statements for $\delta$ follow from $|\sqrt\eta-\sqrt{\eta^0}|\le|\eta-\eta^0|^{1/2}$ and $|\sqrt\eta-\sqrt{\eta^0}|\le|\eta-\eta^0|/\sqrt{\eta^0}$.

\emph{(ii).} The two laws are Gaussian with covariance $\tau^2I_{2J}$ and means differing by $\Phi(\delta_1^2,\vartheta'^0)-\Phi(0,\vartheta'^*)$, $b$ cancelling, so the divergence is $D(\delta_1)^2/2\tau^2$. The choice $\vartheta'=\vartheta'^0+\delta_1^2w$, in $K'$ for small $\delta_1$, gives $\Phi(\delta_1^2,\vartheta'^0)-\Phi(0,\vartheta')=\delta_1^2t_\perp+O(\delta_1^4)$, so $D(\delta_1)\le\delta_1^2\|t_\perp\|+O(\delta_1^4)$. Conversely, let $r_2=\min(r_1,\sigma_{\min}(T')/2C_2)$. If $|\vartheta'-\vartheta'^0|\ge r_2$, Step 1 and the triangle inequality give $\|\Phi(\delta_1^2,\vartheta'^0)-\Phi(0,\vartheta')\|\ge\lambda_1r_2-L\delta_1^2>\delta_1^2\|t_\perp\|+O(\delta_1^4)$ for small $\delta_1$. If $h'=\vartheta'-\vartheta'^0$ has $|h'|<r_2$, the difference is $\delta_1^2t-T'h'+e$ with $\|e\|\le C_2(\delta_1^4+|h'|^2)$; its norm is at least $\sigma_{\min}(T')|h'|/2-\delta_1^2\|t\|-C_2\delta_1^4$, which exceeds the upper bound for $D(\delta_1)$ unless $|h'|=O(\delta_1^2)$, and at least $\|(I-P')(\delta_1^2t-T'h')\|-\|e\|=\delta_1^2\|t_\perp\|-C_2(\delta_1^4+|h'|^2)$. So $D(\delta_1)=\delta_1^2\|t_\perp\|+O(\delta_1^4)$. By Le Cam's two point lemma and Pinsker's inequality, the two error probabilities of any test sum to at least $1-\mathrm{TV}\ge1-(\mathrm{KL}/2)^{1/2}=1-D(\delta_1)/2\tau$, which is $1-c^2/2-O(\tau)$ at $\delta_1=c\,r_\tau^{1/2}$.

\emph{(iii).} Let $\hat\vartheta'$ minimise $\|Y-\Phi(0,\vartheta')\|$ over $K'$. Since $\|Y-\Phi(0,\vartheta'^0)\|\le\|n\|+L\eta^0=O_P(\tau)$, also $\|\Phi(0,\hat\vartheta')-\Phi(0,\vartheta'^0)\|=O_P(\tau)$, and Step 1 with the base point $(0,\vartheta'^0)$ gives $|\hat\vartheta'-\vartheta'^0|=O_P(\tau)$; so with probability tending to one $\hat\vartheta'$ is interior and satisfies the normal equations. Expanding, $Y-\Phi(0,\hat\vartheta')=n+\eta^0t-T'(\hat\vartheta'-\vartheta'^0)+O_P(\tau^2)$ and $T'^{\mathsf T}(Y-\Phi(0,\hat\vartheta'))=O_P(\tau^2)$, so $r=(I-P')(n+\eta^0t)+O_P(\tau^2)$ and $Z=\zeta_\perp+\langle t_\perp,b\rangle/(\tau\|t_\perp\|)+v+O_P(\tau)$, the middle term being at most $\|b\|_2/\tau\to0$. The level and the power follow, and the power is $\tfrac12$ when $v=z_\alpha$, that is $\eta^0=z_\alpha r_\tau$. The least squares fit over $\eta\ge0$ and $\vartheta'$ is consistent in the same way, and in the linearised model the minimum of the residual sum of squares is $\|r\|^2-\max(\langle t_\perp,r\rangle,0)^2/\|t_\perp\|^2+O_P(\tau^3)$; so twice the log likelihood ratio is $\max(Z,0)^2+o_P(1)$, and for $\alpha<\tfrac12$ the events $\max(Z,0)^2>z_\alpha^2$ and $Z>z_\alpha$ coincide.
\end{proof}

Part~(i) is proved, with the conventions of Theorem~\ref{es:thm:concentration}: the noise, a bias bound $\|b\|_2\le B\tau$, the separation and no aliasing conditions, and a rate with no logarithmic factor; the second rate of~(i), $\tau/(\delta_0\|t_\perp\|)$, is claimed only for fixed $\delta_0>0$, since under the spike of Theorem~\ref{es:thm:bf} it fails in the intermediate range. Part~(ii) is proved and shows that no method does better than the rate $\tau^{1/2}$ for the real part, and (iii) is proved for the matched filter.

The scale is visible in the data. We took the window about $60$ at level $0.6$ with taper $\ell=3$ and the samples $\omega_j=2+0.2j$, $j<100$ (the design of the detection experiment below, assumed because the source states none for this window), with a pair at $\tfrac12\pm\delta$ added at $60.089$ to the zeros, the pair fitted with one free common amplitude, and the noise $\tau=1.7\times10^{-11}$ of the source, an estimate from a residual, taken as the standard deviation per real coordinate. In this design $\|t_\perp\|=323.8$ for the free amplitude and $406.5$ for an amplitude fixed by the model, the matched filter has $\delta^*=(1.645\,\tau/\|t_\perp\|)^{1/2}=2.9\times10^{-7}$, and the Bayes factor, which pays an Occam factor, crosses one at about $3\times10^{-7}$ (\texttt{code/es\_bf.py}, \texttt{data/es\_bf\_src.json}, computed, one noise draw per $\delta$: a pair at $10^{-6}$ is already detected, with a profile likelihood gain of $380$ in $2\log$). The data accuracy of the source is $10^{-11}$, and in the design above a real part is resolved to a few times $10^{-7}$. The comparison with Theorem~\ref{es:thm:singular} is with its scale only: the error of the run is deterministic, one noise draw, and Theorem~\ref{es:thm:singular} does not cover the reported runs (Remark~\ref{es:rem:singular}).

\subsection{Bayes factors with the right alternative}

\begin{theorem}[Bayes factors]\label{es:thm:bf}
In the setting of Theorem \ref{es:thm:singular}, put the spike and slab prior on the pair: $\delta=0$ with probability $1-p$, and otherwise $\delta$ with a density $\pi$ continuous and positive on $[0,\delta_+]$. Let $s=\|t_\perp\|/\tau$, $G(Z)=\int_0^\infty u^{-1/2}e^{Zu-u^2/2}\,du$, and $Z$ the statistic of Theorem \ref{es:thm:singular}(iii).
\begin{enumerate}
\item[(a)] If $\eta^0=0$, or $\eta^0=v\,r_\tau$ with $v$ fixed, and $\|b\|_2=o(\tau)$, and $\vartheta'$ has the same prior under the line and under the pair, the Bayes factor for the line against the pair is
\[
\mathrm{BF}_{01}=\frac{2\,s^{1/2}}{\pi(0)\,G(Z)}\,\bigl(1+o_P(1)\bigr);
\]
under the null $Z$ is asymptotically standard Gaussian and $\mathrm{BF}_{01}\asymp\tau^{-1/2}$ in probability.
\item[(b)] Under a pair with fixed $\delta\in(0,\delta_+)$, if $p'$ is bounded below on $K'$ and every point of $K'$ has prior mass at least $c\tau^{D'}$ within distance $\tau$, then $\log\mathrm{BF}_{01}=-D(\delta)^2(1+o_P(1))/2\tau^2$, with $D(\delta)=\delta^2\|t_\perp\|+O(\delta^4)$ as in Theorem \ref{es:thm:singular}(ii).
\item[(c)] Against the alternative of a single zero of width $a_0+\varepsilon$, with $\varepsilon$ having a density positive and continuous at $0$ and $\|(I-P')\partial_\varepsilon\Phi\|>0$, the same data give $\mathrm{BF}_{01}\asymp\tau^{-1}$ in probability under the null.
\end{enumerate}
\end{theorem}

\begin{proof}
{\emergencystretch=8em (a) $\mathrm{BF}_{10}=m_1/m_0$ with $m_1=\int\!\!\int_{\eta>0}e^{\ell(\eta,\vartheta')}\pi_\eta(\eta)\,p'(\vartheta')\,d\eta\,d\vartheta'$ and $m_0=\int e^{\ell(0,\vartheta')}p'(\vartheta')\,d\vartheta'$, $\ell$ the log likelihood ratio to $(0,\vartheta'^0)$. The offset $\Phi(\eta^0,\vartheta'^0)-\Phi(0,\vartheta'^0)=\eta^0t+O(\tau^2)$ is $O(\tau)$ and can be carried in $b$ in Steps 2 and 3 of the proof of Theorem \ref{es:thm:singular}, which then apply to $m_0$ and to $m_1$: the parts outside $U_\tau$ are negligible, and on $U_\tau$ the integrand factorises as in Step 3 with $W$ replaced by $W_0=\langle t_\perp,n+\eta^0t\rangle/(\tau\|t_\perp\|)=Z+o_P(1)$, by the proof of (iii). The $k$ integrals give $p'(\vartheta'^0)I_2(1+o_P(1))$ in both (in $m_1$ it suffices, as in Step 4, to control the $k$ sections for $\eta\le M_\tau r_\tau$, the rest being negligible), and, with $\eta=u\,r_\tau$ and $\varepsilon_\tau/r_\tau\to\infty$,
\[
\int_0^{\varepsilon_\tau}e^{W_0\eta/r_\tau-\eta^2/2r_\tau^2}\,\frac{\pi(\sqrt\eta)}{2\sqrt\eta}\,d\eta=\frac{\pi(0)}2\,r_\tau^{1/2}\,G(W_0)\,(1+o(1)).
\]}
Since $r_\tau^{1/2}=s^{-1/2}$ and $G$ is continuous and positive, $\mathrm{BF}_{10}=\tfrac12\pi(0)s^{-1/2}G(Z)(1+o_P(1))$. Under the null $Z=\zeta_\perp+o_P(1)$, so $G(Z)$ is bounded above and away from $0$ in probability.

(b) Take $\ell$ relative to the truth $(\delta^2,\vartheta'^0)$ and work on $E_\tau$. For the slab, $\ell\le\sup_x(2x\tau\Lambda_\tau-x^2)/2\tau^2=\Lambda_\tau^2/2$ and Step 2 bounds $m_1$ below, so $|\log m_1|=O(\log1/\tau)$. For the spike, $\|u(0,\vartheta')\|\ge D(\delta)$ for every $\vartheta'$, so $\ell(0,\vartheta')\le-(D^2-2D\tau\Lambda_\tau)/2\tau^2$ once $D\ge\tau\Lambda_\tau$; and within distance $\tau$ of the minimiser $\vartheta'^*$, $\|u(0,\vartheta')\|\le D+L\tau$ and $\ell(0,\vartheta')\ge-(D+L\tau)^2/2\tau^2-(D+L\tau)\Lambda_\tau/\tau$, on a set of prior mass at least $c\tau^{D'}$. Hence $\log m_0=-D^2/2\tau^2+O(\Lambda_\tau/\tau)+O(\log1/\tau)$, and $\log\mathrm{BF}_{01}=\log m_0-\log m_1=-D^2(1+o(1))/2\tau^2$.

(c) As in (a), with the regular coordinate $\varepsilon$ in place of $\eta$: the $\varepsilon$ integral is $\pi_\varepsilon(0)\,\tau\,\|(I-P')\partial_\varepsilon\Phi\|^{-1}\int e^{Z_\varepsilon v-v^2/2}\,dv\,(1+o_P(1))$, over the half line or the line allowed by the sign of $\varepsilon$, which is of order $\tau$; so $\mathrm{BF}_{10}\asymp\tau$.
\end{proof}

The formula needs the statistic $Z$ and the scale $s$, so that its use on data is a computation, and the one for the real zero at $201.264752$ is an illustration under assumptions. The source does not give the noise level of that window. With $\pi$ uniform on $(0,\tfrac12)$, the design of the detection experiment below (level $0.6$, $\ell=3$, $\omega_j=2+0.2j$, an assumption) and $\tau=1.7\times10^{-11}$ assumed, the simple zero, of amplitude $A^0=\tfrac12$, has $\|t_\perp\|=148.6$ and $s=\|t_\perp\|/\tau=8.74\times10^{12}$; the formula of Theorem~\ref{es:thm:bf} gives $\log_{10}\mathrm{BF}_{01}=6.14$ at $Z=0$ and $5.51$ at $Z=1.7$, and the integral computed exactly gives a median $6.10$ over $40$ noise draws, with range $5.55$ to $6.48$. The single zero alternative gives $11.83$ at $Z_1=0$ and $10.92$ at $Z_1=1.7$, and $9.66$ to $11.78$ for the other ten central zeros at $Z_1=0$, the spread arising from the norm of the projected tangent, which varies from $0.2$ to $29$ over the central zeros (computed, \texttt{code/es\_bf.py}, \texttt{data/es\_bf\_src.json}). So, under these assumptions, one window of data at this accuracy is worth a factor of about $10^6$ per zero in favour of the line against the pair, against about $10^{12}$ against the single zero alternative: the functional equation takes the square root of the evidence (computed in the free amplitude model; heuristic beyond it).

\section{An adaptive prior, and a ReLU network}\label{es:sec:adaptive}

\subsection{Where to look: the ordinates}

The prior of Theorem~\ref{es:thm:concentration} put the ordinates uniformly on the window. A prior that regularises where the zeros are has two pieces of information to use. The first is unconditional: the zeros have density $d(t)=\log(t/2\pi)/2\pi$ at height $t$, with an explicit error, so the number in a window of fixed length $W$ is $Wd(c)+O(\log c)$ (proved, Riemann and von Mangoldt). The second is conjectural and enters only as a prior: conjecturally (the pair correlation theorem of Montgomery \cite{Montgomery73}, under RH and for restricted support, and the computations of Odlyzko \cite{Odlyzko87}) the local statistics of the ordinates are those of the sine kernel determinantal process. So take the prior on the ordinates in a window to be the determinantal point process with kernel
\[
K(x,y)=\frac{\sin\pi d(c)(x-y)}{\pi(x-y)},
\]
which adapts to the height through $d(c)$ and repels: its Janossy density vanishes like $\prod_{i<j}(x_i-x_j)^2$ at coincidences. Three consequences. At a configuration of distinct ordinates the rates of Theorems~\ref{es:thm:concentration} and~\ref{es:thm:singular} are unchanged, the prior density being continuous and positive there, since the prior mass of a ball of radius $\tau$ about it is still of order $\tau^{\dim}$; only the constants move. The posterior on the number $M$ of zeros in the window is pulled toward $Wd(c)$, a prior tendency and the Bayesian form of the count that Section~\ref{es:sec:above} could not certify at $3\times10^{12}$ (heuristic). And spurious components that sit on top of genuine ones are penalised, a priori and by the two point correlation, by the factor $(\pi d\,\lvert x_i-x_j\rvert)^2/3$, which is $1\times10^{-4}$ at the density $d=0.569$ of height $224$ and a separation $0.0097$, and which may help the many window scheme remove them (heuristic). The repulsion is a conjecture about the zeros and enters only as a prior. A determinantal process has almost surely no two points at one ordinate, so the process is put on the distinct ordinates, at the intensity $d(c)$, a pair at $\tfrac12\pm\delta$ counting as one point marked by $\delta$; put on the zeros themselves it would give prior probability zero to every off line pair, and so decide the real parts before the data. With this convention the ordinate prior is the same under the line and under the pair, and to leading order as $\tau\to0$ it leaves the Bayes factor of Theorem~\ref{es:thm:bf} unchanged (Remark~\ref{es:rem:dpp}); at finite $\tau$ it can still move the posterior of the real parts, as any prior on parameters that the likelihood couples with them can.

\begin{remark}\label{es:rem:dpp}
For a determinantal process with continuous kernel $K$ and intensity $d$ the second correlation function is $\rho_2(x,y)=d^2-K(x,y)^2\le d^2$, so the expected number of ordered pairs of distinct points of a window of length $W$ at distance less than $\epsilon$ is at most $2\epsilon Wd^2$, which tends to $0$ with $\epsilon$: almost surely no two points coincide, while a pair at $\tfrac12\pm\delta$, and a double zero, put two zeros at one ordinate. On the distinct ordinates the prior density $p'$ of $\vartheta'$ in Theorem~\ref{es:thm:singular} (the Janossy density of the ordinates times the priors of the amplitudes and widths) does not depend on $\eta$, and it is continuous and positive at distinct ordinates; in the proof of Theorem~\ref{es:thm:bf}(a) the integrals over $\vartheta'$ give $p'(\vartheta'^0)I_2(1+o_P(1))$ in both marginal likelihoods, so $p'$ cancels from $\mathrm{BF}_{01}$ to leading order. At finite noise it does not cancel from the posterior of a real part: if $Y_1=\gamma+a+\sigma\epsilon_1$ and $Y_2=\gamma+\sigma\epsilon_2$, with $\epsilon_1,\epsilon_2$ independent standard normal and independent priors $\gamma\sim N(m,v_\gamma)$, $a\sim N(0,v_a)$, then
\[
E[a\mid Y_1,Y_2]=\frac{v_a(Y_1-m_2)}{v_a+v_2+\sigma^2},\qquad m_2=\frac{\sigma^2m+v_\gamma Y_2}{v_\gamma+\sigma^2},\quad v_2=\frac{v_\gamma\sigma^2}{v_\gamma+\sigma^2},
\]
which depends on the prior mean $m$ of the ordinate whenever $\sigma>0$ (given $Y_2$, $\gamma$ is $N(m_2,v_2)$, and then $Y_1$ given $a$ is $N(m_2+a,v_2+\sigma^2)$).
\end{remark}

\subsection{What to believe: the real parts}

For the real parts the adaptive prior is the spike and slab on pairs of Theorem~\ref{es:thm:bf}: the functional equation supplies the structure, the spike at $\delta=0$ supplies the shrinkage, and the slab keeps the alternative honest. This is the prior under which a computation like that of Section~\ref{es:sec:above} should be read, and Theorem~\ref{es:thm:bf} says what it can deliver: under the line, Bayes factors of order $\tau^{-1/2}$ in probability per pair, the other pairs held on the line, in the free amplitude model (the posterior odds are $(1-p)\,\mathrm{BF}_{01}/p$), and a resolution of $(\tau/\|t_\perp\|)^{1/2}$ in the real part.

\subsection{A ReLU network}

The question was whether a ReLU network can serve as the adaptive prior or regulariser. Polson and Ro\v{c}kov\'a \cite{PolsonRockova18} proved that deep ReLU networks with a spike and slab prior on their weights have posteriors that contract at the near minimax rate for H\"older functions, adapting to the unknown smoothness. That theorem is about nonparametric regression and gives no rate for the window model, and what it suggests here is heuristic; part~(ii) of Theorem~\ref{es:thm:singular} bounds every method, networks included. Where a network can help is in amortisation: it learns a classifier of the line against the pair, which at the population optimum is the posterior probability of the pair under the training law, and costs one pencil fit and one forward pass per window. In the author's experiment the training law is a recentred circular unitary ensemble of $40$ points truncated to the window, with a log uniform slab on $[10^{-7.5},0.09]$ for the displacement of the pair, which is neither the sine kernel process nor the spike and slab of Theorem~\ref{es:thm:bf} (this is read from the author's script). That experiment, with two ReLU detectors (a network with two hidden layers of $256$ units on the raw transform and the same network on the residual of the fit on the line), was run by the author and is unchecked: its two rows were not reproduced here and are not printed.

The pencil and a plug in matched filter were rerun for this book on the window model of Theorem~\ref{es:thm:taper} at level $0.6$ with taper $\ell=3$ and $100$ samples $\omega_j=2+0.2j$ (the sampling that the author's script appears to use; the chapter does not give it), ordinates drawn from the circular unitary ensemble scaled to the density at height $200$, and the zero nearest the centre either a double zero on the line or a pair at $\tfrac12\pm\delta$; the noise has complex root mean square $\tau$, that is standard deviation $\tau/\sqrt2$ per real coordinate. Two detectors, each calibrated to five per cent false alarms on $1200$ null windows: the pencil statistic $\max\lvert\hat a-a_0\rvert$ over the components within $0.5$ of the centre, and a plug in version of the matched filter of Theorem~\ref{es:thm:singular}(iii) in which the tangent $t$ is that of the central zero at the pencil's ordinate and $t_\perp$ is orthogonal to the tangent columns (amplitude and pole) of every pencil component (this is our reading of ``built from the pencil's components''; the source does not define it, and no efficiency claim is made for any plug in filter). The quantity $\delta^*$ is the smallest displacement detected with power one half.

\begin{center}
\footnotesize
\setlength{\tabcolsep}{3pt}
\begin{tabular}{lccccc}
\toprule
& \multicolumn{5}{c}{detectable $\delta^*$ at noise $\tau=$}\\
detector & $10^{-10}$ & $10^{-8}$ & $10^{-6}$ & $10^{-4}$ & $10^{-2}$\\
\midrule
benchmark $0.0537\,\tau^{1/2}$ & $5.4\times10^{-7}$ & $5.4\times10^{-6}$ & $5.4\times10^{-5}$ & $5.4\times10^{-4}$ & $5.4\times10^{-3}$\\
pencil & $8.5\times10^{-7}$ & $7.0\times10^{-6}$ & $6.1\times10^{-5}$ & $9.1\times10^{-4}$ & $7.2\times10^{-3}$\\
plug in filter & $1.0\times10^{-6}$ & $1.3\times10^{-5}$ & $1.3\times10^{-4}$ & $1.0\times10^{-3}$ & $1.2\times10^{-2}$\\
\bottomrule
\end{tabular}
\end{center}

\noindent The benchmark is Theorem~\ref{es:thm:singular}(iii) at the half power point, $\delta^*=(1.645\,\tau'/\|t_\perp\|)^{1/2}$, with the median $403$ of $\|t_\perp\|$ over $300$ windows ($10$th and $90$th percentiles $382$ and $407$) and $\tau'=\tau/\sqrt2$, the standard deviation per real coordinate; with $\tau$ itself as the standard deviation per real coordinate the constant is $0.0639$ and the row is $6.4\times10^{-7}$ to $6.4\times10^{-3}$. It is the half power point at the median, a benchmark and not a bound (a law of $\|t_\perp\|$ with two values can have a smaller average power at it). The pencil and plug in rows are the means of two seeds ($150$ windows per $\delta$, $16$ values of $\delta$; \texttt{code/es\_detect.py}, \texttt{code/es\_detect\_all2.py}, \texttt{data/es\_detect\_src.log}, computed).

Two findings. First, the law $\delta^*\propto\tau^{1/2}$ of Theorem~\ref{es:thm:singular} holds over eight decades of noise for the pencil, with $\delta^*/\tau^{1/2}$ between $0.057$ and $0.093$, within a factor $1.7$ of the benchmark; a data accuracy of $10^{-10}$ resolves a real part to about $10^{-6}$ in this design. Second, the plug in filter is $1.1$ to $2.1$ times worse than the pencil in our implementation. The source's rows were not reproduced; they depend on the pencil's order rule and on the definition of the plug in filter, which the source does not give.

So the answer to the question has two parts. A ReLU network is not a regulariser that changes what can be found: the bound in part~(ii) of Theorem~\ref{es:thm:singular} holds for it, and a network trained on residuals is a learned test that cannot beat it. Whether such a learned test is more robust than the pencil is a matter for the author's experiment, which is unchecked. The regularisation that matters is in the prior: the sine kernel for where the zeros are, and the spike and slab on pairs for what their real parts may be. None of this reaches the hypothesis. Each window at accuracy $\tau$ gives, under the line, Bayes factors of order $\tau^{-1/2}$ in probability per pair, the other pairs held on the line, in the free amplitude model, and a resolution of order $\tau^{1/2}$ in the real part, finitely many zeros at finite precision, as before.

\section{What this does and does not do}

Sparse methods solve the deconvolution. The problem $P_{1/2}*\mu=u_1$ is exponentially ill posed for a general $\mu$ and polynomially ill posed for a sum of atoms, and the matrix pencil, which works in the second model, recovers the atoms to ten digits where EM, which works in the first, reaches three. The structural reason is Theorem~\ref{es:thm:taper}: the explicit formula, read as a deconvolution and tapered, is a sum of damped exponentials, Prony's problem, and the level one data determine the poles. Two further things follow. The pencil does not need to assume a width: it returns $a_\rho$ and so places each zero in the strip, which the backward flow at a fixed level cannot do, since that flow presupposes a width. And the certificate question of Chapter~\ref{chap:cx} changes character: a certified version of the pencil would need interval versions of the singular value decomposition and the eigenvalue problem, which exist, and a bound linking the matrix model to $\xi$ (the remainder of Theorem~\ref{es:thm:taper} and the far zeros), which does not, and would then give certified zeros from the line $\Re s=1$; we have not done it, and it is the first open problem below.

What it does not do is also clear. The data $u_1$ to twenty digits are $\Re(\xi'/\xi)$ on $\Re s=1$ to twenty digits, which is computed from $\zeta$ by the same algorithms that find zeros directly; the pencil is a different route from the same information, not a cheaper one. It reads the real parts; it does not constrain them. The one thing it offers the Thorin programme is a sharper statement of what the programme computes: the backward flow at level $\theta$ is the resolvent trace $\pi^{-1}\Re(\xi'/\xi)(\theta+iv)$, whose Fourier transform is $\sum_\rho\operatorname{sgn}(\theta-\sigma_\rho)e^{-\lvert\theta-\sigma_\rho\rvert\lvert\omega\rvert}e^{-i\omega\gamma_\rho}$ (Proposition~\ref{es:prop:level}), and the Riemann hypothesis is the statement that every damping rate $\lvert\theta-\sigma_\rho\rvert$ at level $\theta$ is $\theta-\tfrac12$, which is what the pencil measures and what no measurement proves.

A Bayesian reader will recognise the two models as two priors on the same unknown; we state the comparison as an analogy, since no posterior was computed here. The grid prior of Chapter~\ref{chap:cx}, a product of gamma laws on the cells, is a prior on densities, and the EM solution is a point estimate under that model, with its $10^{-3}$. A spike prior, a point process of atoms with amplitudes bounded away from zero, is a prior on sums of atoms; in the setting of Theorem~\ref{es:thm:concentration} its posterior concentrates at the true atoms at the rate of the noise, and the matrix pencil, which is exact up to the remainder of Theorem~\ref{es:thm:taper} when the data come from finitely many atoms, reaches the same atoms without computing a posterior. The deconvolution is the same; what changed is the model, and in an exponentially ill posed problem the model is everything.

\section{Questions taken up}\label{es:sec:worked}

This section takes up the questions left open by the first sections: how to certify the pencil's output, how close pairs behave, what the levels above one half add, which level to use, how to combine windows, and whether the method works at a large height. The tables of this section are the authors' computations at heights near $200$, not rerun for this book, except where a sentence says otherwise.

\subsection{Certifying the pencil}\label{es:sec:certpencil}

The pencil proposes; ball arithmetic disposes. For each of the eleven zeros within $10$ of the centre $c=200$, with the pencil's value $\gamma$, two checks were made in Arb at $80$ bits. First, $\xi(\tfrac12+it)$ changes sign on $[\gamma-2\times10^{-9},\gamma+2\times10^{-9}]$, which certifies a zero on the line within $2\times10^{-9}$ of $\gamma$. Second, the winding number of $\xi$ around the circle of radius $10^{-3}$ about $\tfrac12+i\gamma$, computed as $(2\pi i)^{-1}\oint\xi'/\xi$ by Arb's rigorous integrator, is the ball $1.0000\pm10^{-9}$, hence exactly $1$: there is one zero in the disc, it is simple, and it is the one on the line. All eleven pass. This is not the same as Rump's verified eigenvalue bounds \cite{Rump10} for the pencil itself, which would enclose the eigenvalues of the matrix model and would still need a bound linking that model to $\xi$; it is what one wants in practice: the pencil locates, the argument principle certifies, and together they give certified zeros from level one data in the time the pencil takes plus a few seconds per zero. For a zero off the line the same winding number around a circle of radius less than $\lvert\sigma-\tfrac12\rvert$ about the pencil's $\sigma+i\gamma$ would certify it with the same rigour, which is more than a sign change can do.

\subsection{Close pairs}

The pencil's resolution of a close pair was tested on synthetic data built from Theorem~\ref{es:thm:taper}: the Lorentzians of the true zeros within $15$ of the centre, with those within $1.5$ of the centre removed and a pair at $c\pm d/2$ put in their place, for $d$ from $1$ down to $0.01$, with Gaussian noise $\eta=10^{-16}$ or $10^{-12}$ on the samples and bandwidths $\Omega=25$, $45$, $65$. The error in the recovered separation is below $10^{-12}$ at $\eta=10^{-16}$ and below $5\times10^{-9}$ at $\eta=10^{-12}$, for every $d$ down to $0.01$ and every $\Omega$: the loss from closeness is a factor of a few hundred at most over two decades of separation, a polynomial loss consistent with the bounds of Donoho and of Li and Liao \cite{Donoho92,LiLiao21}, not the exponential one of Proposition~\ref{es:prop:resolution}. The closest real pair below height $1420$ reported by the authors, at $1329.043518$ and $1329.205019$ with gap $0.1615$ (not checked here), is recovered from the real data to about $2\times10^{-8}$ at level $1$ and to about $3\times10^{-11}$ at level $0.6$. Down to a separation of $0.01$, close pairs are not the limitation of the method.

\subsection{Levels above one half}

The pencil was run on $u_\theta(v)=\pi^{-1}\Re(\xi'/\xi)(\theta+iv)$ for $\theta=1$, $\tfrac34$ and $0.6$ on the same window. The zeros enter with a sign, and the sign matters for a zero to the right of $\theta$.

\begin{proposition}[levels below the zeros]\label{es:prop:level}
Let $\tfrac12<\theta\le1$ with $\theta\ne\sigma_\rho$ for every $\rho$, and $u_\theta(v)=\pi^{-1}\Re(\xi'/\xi)(\theta+iv)$. Then
\[
\begin{gathered}u_\theta(v)=\sum_\rho\operatorname{sgn}(\theta-\sigma_\rho)\,P_{|\theta-\sigma_\rho|}(v-\gamma_\rho),\\ \widehat{u_\theta}(\omega)=\sum_\rho\operatorname{sgn}(\theta-\sigma_\rho)\,e^{-|\theta-\sigma_\rho||\omega|}e^{-i\omega\gamma_\rho},\end{gathered}
\]
and for $\omega\ell^2>1$ the tapered transform $F_\theta(\omega)=\int u_\theta(v)w(v)e^{i\omega(v-c)}dv$ is
\[
F_\theta(\omega)=\sum_\rho\operatorname{sgn}(\theta-\sigma_\rho)\,w(\gamma_\rho+i|\theta-\sigma_\rho|)\,e^{-|\theta-\sigma_\rho|\omega}e^{i\omega(\gamma_\rho-c)}+R_\theta(\omega),
\]
with $R_\theta$ bounded as in Theorem~\ref{es:thm:taper} with $a_\rho$ replaced by $|\theta-\sigma_\rho|$. A zero with $\sigma_\rho>\theta$ therefore gives the pole $e^{\Delta(-(\sigma_\rho-\theta)+i(\gamma_\rho-c))}$, inside the unit circle on the ring of a zero at $2\theta-\sigma_\rho$, with an amplitude of sign opposite to the taper profile, while its partner $1-\sigma_\rho<\theta$ gives a pole at the same ordinate with the profile's sign. The Riemann hypothesis holds if and only if every damping rate $|\theta-\sigma_\rho|$ equals $\theta-\tfrac12$.
\end{proposition}

\begin{proof}
With the zeros paired, $\xi'/\xi(s)=\sum_\rho(s-\rho)^{-1}$, and $\Re(\theta+iv-\rho)^{-1}=(\theta-\sigma_\rho)/((\theta-\sigma_\rho)^2+(v-\gamma_\rho)^2)=\pi\operatorname{sgn}(\theta-\sigma_\rho)P_{|\theta-\sigma_\rho|}(v-\gamma_\rho)$; the series converges absolutely, its terms being $O(\gamma_\rho^{-2})$. Transform termwise with $\widehat{P_b}(\omega)=e^{-b|\omega|}$. For $F_\theta$ apply the proof of Theorem~\ref{es:thm:taper} to each $P_{|\theta-\sigma_\rho|}$, whose width lies in $(0,1)$, below $\omega\ell^2$. A pole of modulus $e^{-\Delta b}$ is read as the width $b$, that is as the real part $\theta-b$, which for $b=\sigma_\rho-\theta$ is $2\theta-\sigma_\rho$. For the last statement: if $|\theta-\sigma_\rho|=\theta-\tfrac12$ for every $\rho$, each $\sigma_\rho$ is $\tfrac12$ or $2\theta-\tfrac12$; but if $\sigma_\rho=2\theta-\tfrac12$ then $1-\sigma_\rho=\tfrac32-2\theta$ is also the real part of a zero, with rate $3(\theta-\tfrac12)\ne\theta-\tfrac12$.
\end{proof}

For $\sigma_\rho<\theta$ the poles are $e^{\Delta(-(\theta-\sigma_\rho)+i(\gamma_\rho-c))}$ (equation~\eqref{es:eq:lorentz} with $a_\rho=\theta-\sigma_\rho$ and Proposition~\ref{es:prop:level}). Three findings, from the authors' run.

\begin{center}
\small
\begin{tabular}{lcccc}
\toprule
$\theta$ & position error, max & median & error in $\sigma_\rho$, max & median\\
\midrule
$1$ & $4.5\times10^{-10}$ & $1.8\times10^{-10}$ & $7.9\times10^{-10}$ & $2.3\times10^{-10}$\\
$3/4$ & $3.3\times10^{-11}$ & $1.5\times10^{-11}$ & $5.3\times10^{-11}$ & $1.4\times10^{-11}$\\
$0.6$ & $2.4\times10^{-11}$ & $8.0\times10^{-12}$ & $2.3\times10^{-11}$ & $7.2\times10^{-12}$\\
\bottomrule
\end{tabular}
\end{center}

\noindent First, the real parts agree across levels: $\sigma_\rho=\theta-a_\rho$ is $0.50000000$ for every one of the eleven central zeros at all three levels, to eight digits, a consistency check that the computation of $u_\theta$ passes. Second, lower levels are better conditioned, by an order of magnitude from $1$ to $\tfrac34$ and a little more to $0.6$: the damping $e^{-(\theta-1/2)\omega}$ is milder, so the signal at the top of the band is $e^{-4.5}$ at $\theta=0.6$ against $e^{-22.5}$ at $\theta=1$, and the pencil's noise amplification falls accordingly. Third, and this is the point of the exercise, nothing in the method needs the level to be above the zeros: $u_\theta$ is defined for every $\theta$ that is not the real part of a zero, and by Proposition~\ref{es:prop:level} the Lorentzians of zeros to the right of $\theta$ have negative width and appear as poles inside the unit circle, of modulus $e^{-\Delta(\sigma_\rho-\theta)}$, with amplitudes of the opposite sign to the taper profile, so that the pencil reads them at $2\theta-\sigma_\rho$; the sign of the amplitude is the test, a test that the amplitude filter used in these runs, which compares moduli, does not perform. The pencil at level $\theta$ could therefore be a test for zeros in $\Re s>\theta$ on a window through the sign of the amplitude (not implemented in the runs reported here, open), with a certificate available by the winding number of Section~\ref{es:sec:certpencil}. The recovered poles at the three levels lie on three concentric rings, one per level, all at $\sigma_\rho=\tfrac12$: the argument of a pole is the ordinate and its modulus is $e^{-\Delta(\theta-\sigma_\rho)}$. The selection of components in these runs compares the amplitude with the taper profile at the width $\theta-\tfrac12$ of a zero on the line (two independent readings of the authors' scripts agree), so that these runs presume the real part $\tfrac12$ in the selection step.

\subsection{The level as a cost parameter}

Which level should one use? In a separate run, the pencil was run at $\theta=1$, $\tfrac34$, $0.6$, $0.55$, $0.52$ and $0.51$ on the window about $200$, with the grid spacing reduced as the level drops so that the leading aliasing error $e^{-(\theta-1/2)(2\pi/h-\Omega)}$ of the transform stays below $3\times10^{-12}$: $h=0.02$ down to $0.6$, then $0.01$, $0.004$ and $0.002$.

\begin{center}
\small
\begin{tabular}{lcccc}
\toprule
$\theta$ & $h$ & data points & position error, median & error in $\sigma_\rho$, median\\
\midrule
$1$ & $0.02$ & $5001$ & $1.8\times10^{-10}$ & $2.3\times10^{-10}$\\
$3/4$ & $0.02$ & $5001$ & $1.5\times10^{-11}$ & $1.3\times10^{-11}$\\
$0.6$ & $0.02$ & $5001$ & $8.0\times10^{-12}$ & $7.2\times10^{-12}$\\
$0.55$ & $0.01$ & $10001$ & $3.7\times10^{-12}$ & $4.2\times10^{-12}$\\
$0.52$ & $0.004$ & $25001$ & $2.8\times10^{-12}$ & $3.9\times10^{-12}$\\
$0.51$ & $0.002$ & $50001$ & $9.9\times10^{-13}$ & $1.4\times10^{-12}$\\
\bottomrule
\end{tabular}
\end{center}

\noindent The error falls by a factor twenty from level $1$ to level $0.6$ and then only slowly, by another factor eight down to $0.51$, while the number of data points, and with it the cost of the data, grows like $1/(\theta-\tfrac12)$. The gain presumably saturates because the pencil's error is set by the lower frequencies, where the signal is large at every level, and not by the top of the band where the damping $e^{-(\theta-1/2)\Omega}$ differs most (heuristic). So the answer to the problem is that the level is a weak parameter: anything in $[0.55,0.75]$ is within a factor of twenty of the best, and $\theta=0.6$ with the grid $0.02$ is a good default, costing the same data as level $1$ and giving twenty times the accuracy. There is no level below which the data become as expensive as finding zeros directly, since Arb evaluates $\zeta$ at a point of the strip at about the cost of a point on the line at the same height; what grows is the number of points, and only linearly.

\subsection{Many windows}

A range of heights is covered by combining windows, and the simplest rule works: tapers of width $\ell=10$ centred every $10$ units, each zero taken from the window whose centre is nearest to it, with half a unit of overlap. On $[100,400]$ at level $0.6$ this is $31$ windows, and all $173$ zeros in the range are recovered (the count agrees with $N(400)-N(100)=173$), with median position error $2.4\times10^{-12}$, maximum $2.6\times10^{-11}$, and similar figures for the real parts; the pencils take four seconds, the data on $[40,460]$ at spacing $0.02$ seven. Seven spurious components survive the amplitude filter, all at the edges of windows and all with amplitude at the lower limit of the filter; a stricter filter, or the agreement of two overlapping windows, should remove them. The cost to height $T$ is about $T/h$ evaluations of $\xi'/\xi$ off the line, fifty per unit height at $h=0.02$, each about as expensive as an evaluation of $\zeta$ at that height, plus $T/10$ pencils whose order grows like $\log T$ with the density of zeros; a sign change scan needs a few evaluations of $Z$ per zero, on the line. Both counts grow like $T$ up to logarithms and each evaluation costs about the same, so the constants decide, and they favour the direct methods by a wide margin; the pencil's advantage is the real parts, which the sign changes never see.

\subsection{A window near the limit of the verified range}\label{es:sec:above}

We ran the pencil at a large height, near the limit of the rigorous verification: Platt and Trudgian \cite{PT21} verify the hypothesis for $0<\gamma\le3\,000\,175\,332\,800$, and we took the window centred at $c=3\times10^{12}+100$, inside that range, at level $\theta=0.6$ with taper $\ell=2$, chosen because the zeros there are $4.3$ per unit height and a window of $\pm6$ holds about fifty of them. No prime sum is used: the data are values of $\xi'/\xi$ at $\Re s=0.6$. They are $u_{0.6}$ at $1001$ points at spacing $0.02$, each from the first two Taylor coefficients of $\xi$ at $0.6+iv$ in Arb at $160$ bits, $3.8$ seconds a point on our machine, since the Riemann and Siegel sum has $7\times10^5$ terms there and the Gamma factor costs $46$ bits of the working precision at this height; the $1001$ points took $32$ minutes on two cores. The pencil used $\omega\in[5,85]$ at spacing $0.2$, a Hankel matrix of order $200$, and model order $90$ by the singular value gap. This is the authors' computation (scripts \texttt{pencil\_3e12.py} and \texttt{pencil\_3e12\_data.py} of the source, which this book did not run), computed and not rerun.

\begin{resultpar}\label{es:res:window3e12}It returned $52$ components with amplitudes on the taper's profile within $6$ of the centre, at mean spacing $0.2325$ between the extreme ordinates, the density $\log(c/2\pi)/2\pi=4.28$ per unit predicting $51$ (the smooth count is $51.36$). Every one of them has
\[
\sigma_\rho=0.5000000\pm 1.4\times10^{-8},
\]
the real part read from the modulus of the pole, which the fit does not constrain. The components are selected by comparing the moduli of their amplitudes with the taper profile computed at the width $\theta-\tfrac12$ of a zero on the line, so a zero far from the line would have been discarded by the selection; that the selection discarded no zero on the line is shown by the count of sign changes below. The ordinates are then certified on the line by sign changes of Hardy's function $Z(t)=e^{i\vartheta(t)}\zeta(\tfrac12+it)$, computed in balls with $\vartheta$ from $\log\Gamma$ at $160$ bits, on the intervals $[\gamma-10^{-7},\gamma+10^{-7}]$: all $52$ of the $52$ intervals contain a sign change, which places a zero on the line within $10^{-7}$ of each pencil ordinate (the authors' certificate, not rerun). An independent scan of Hardy's function in ball arithmetic at $160$ bits (\texttt{code/ex\_zeros3e12.py}, \texttt{data/ex\_zeros3e12.json}) finds exactly $52$ certified sign changes in $[c-6,c+6]$, at ordinates from $c-5.91$ to $c+5.95$, with brackets of width below $10^{-10}$; the comparison with the pencil's ordinates, one by one, has not been made (the pencil's list is not available), and a sign scan cannot see two zeros inside one step, so the count $52$ is not a proof of completeness. The winding numbers, which closed the certificate for the eleven zeros near height $200$, did not complete here. Arb's rigorous integrator, fed $\xi'/\xi$ on a circle of radius $0.03$ about the pencil's zero nearest the centre, ran for over an hour at $3.8$ seconds an evaluation without converging, and the alternative of enclosing $\xi$ on arcs of the circle by ball evaluation failed because Arb's enclosure of $\zeta$ at a ball argument of radius $10^{-4}$ at this height already contains zero: the derivative bounds it uses are not tuned to the Riemann and Siegel range. So what is certified at $3\times10^{12}$ is the existence of a zero on the line within $10^{-7}$ of each of the $52$ ordinates the pencil read, and what is not certified is that these are all the zeros in the window, which would follow from a winding number or from Turing's method \cite{Turing53} with the explicit bounds for $S(T)$ \cite{Trudgian11}, and which the pencil's amplitude profile supports without proving. A rigorous argument principle at large height needs certified nonzero values of $\zeta$ on a mesh of the circle together with an explicit bound for $|\zeta'|$ on a small disc, available from the explicit convexity bounds, and is the natural next step.\end{resultpar}

What this is and is not. It is a demonstration that the pencil runs at a height near the limit of the rigorous verification, recovers the zeros of a window from data on the line $\Re s=0.6$, and reports their real parts, which the fit does not constrain, although the selection of components presumes the real part $\tfrac12$. It is not a finding about the zeros, and the window lies inside the verified range. Zeros at and far above this height have been computed before: Gourdon's unrigorous computation of the first $10^{13}$ zeros \cite{Gourdon04} reaches about $2.4\times10^{12}$, and Odlyzko's computations near the $10^{20}$th and higher zeros \cite{Odlyzko92,Odlyzko01} located long runs of zeros on the line by the same sign change test used here (these heights and attributions are the source's and were not checked here). Nor is reading a real part a stronger test than counting: an off line pair in a window shows up in Turing's method as two missing sign changes against the count from $N(T)$, and the pencil adds nothing to that test. Two of its outputs are also weaker than they look. The real parts $0.5000000\pm1.4\times10^{-8}$ come from a floating point pencil and are an estimate, not a certificate; what is certified is a sign change of $Z$ near each ordinate, which is the standard verification. And since the count in the window is not certified, an off line pair the pencil missed is not excluded by anything computed here. What the run does establish is narrower: the deconvolution route of this chapter, from data on $\Re s=0.6$ to positions and real parts, works at height $3\times10^{12}$, with real parts within $1.4\times10^{-8}$ of $\tfrac12$, at a cost of half an hour per window, and its output agrees with the authors' sign change certificate for every zero it reports.

\subsection{Above the verified range}\label{es:sec:above13}

The verified range of Platt and Trudgian ends at $3\,000\,175\,332\,800$, and the window centred at $c=10^{13}+100$ lies above it, by a factor $3.33$ in height. We repeated the computation there with scripts written for this book from the mathematics (the scripts of the source were not used): \texttt{code/es\_pencil.py} generates the data and runs the pencil, and \texttt{code/es\_zeros1e13.py} is the independent scan of Hardy's function; the results are in \texttt{data/es\_pencil1e13.json}, \texttt{data/es\_zeros1e13.json} and, for the data, \texttt{data/es\_u06\_1e13.json}. The setting is that of the previous subsection: level $0.6$, taper $\ell=2$, $1001$ values $u_{0.6}(v_k)$ at $v_k=c-10+0.02k$, each a ball at $160$ bits from the first two Taylor coefficients of $\zeta$ at $0.6+iv_k$ (the largest radius is $1.6\times10^{-27}$, and re-evaluating the first point of every task at $200$ bits changes it by at most $1.7\times10^{-47}$), at $2.6$ seconds a point on one core, $2634$ core seconds in all; the pencil uses $\omega\in[5,85]$ at spacing $0.2$ ($401$ samples) and a Hankel matrix of order $200$ in double precision, in $0.7$ seconds. The model order is $132$, the number of singular values above ten times the plateau of the spectrum: the spectrum decays gradually, a gap rule is not usable, and a smaller order loses zeros.

\begin{resultpar}\label{es:res:pencil1e13}The pencil returns $53$ components with amplitude on the taper's profile within $6$ of the centre, at ordinates from $c-5.7$ to $c+5.8$, the smooth count of the window being $53.66$; the ratio of the amplitude to the profile is $1$ to within $1.6\times10^{-7}$ for all $53$. The real part of each pole is read from its modulus, and the filter that selects the components compares the amplitude with the taper profile computed at the pole's own width, so that no real part is assumed (this is the difference from the runs of the source, whose selection uses the width $\theta-\tfrac12$). The largest $\lvert\sigma_\rho-\tfrac12\rvert$ over the $53$ poles is $4.7\times10^{-9}$ (median $7.1\times10^{-11}$), an estimate from a double precision pencil and not a certificate.\end{resultpar}

\begin{resultpar}\label{es:res:zeros1e13}The scan finds $53$ certified sign changes of Hardy's function in $[c-6,c+6]$: each step of $0.02$ has opposite signs of $Z$ at its ends in ball arithmetic at $160$ bits, and the bisection brackets of width $10^{-10}$ were certified at both ends again at $192$ bits. Every pencil ordinate lies within $10^{-7}$ of the midpoint of a certified bracket, and every certified bracket has a pencil ordinate: all $53$ are matched, no pencil component and no sign change is left over, and the largest ordinate difference is $3.9\times10^{-9}$ (median $5.7\times10^{-11}$). The result is stable: with the model order changed by $\pm5$, $53$ components remain, the ordinates move by at most $3.1\times10^{-9}$ and the real parts by at most $2.1\times10^{-9}$. Against the scan, the largest ordinate difference over the model order is $1.1\times10^{-3}$ at order $70$, $1.9\times10^{-4}$ at $80$, $4.8\times10^{-6}$ at $90$, $5.4\times10^{-9}$ at $95$ and $4.9\times10^{-9}$ at $100$, and below order $65$ the pencil finds only between $8$ and $49$ of the $53$ zeros. The same pencil was validated on known zeros, at $c=200$ (seven zeros) and $c=1000$ (ten zeros): all are matched, the largest ordinate difference is $7.1\times10^{-15}$ and $1.4\times10^{-14}$, and the largest $\lvert\sigma-\tfrac12\rvert$ is $1.7\times10^{-14}$ and $2.3\times10^{-14}$.

The status is this. Computed: the $53$ components, their ordinates and real parts. Certified: $53$ zeros of $\zeta$ lie on the critical line in $[c-6,c+6]$, each located to $10^{-10}$ (this is a certified computation above the height verified in the literature, but it certifies only these $53$ zeros on the line). Not certified: the count (a sign scan cannot see two zeros inside one step of $0.02$, the smallest gap between consecutive zeros found is $0.0326$, so a pair closer than the step is not excluded, and no Turing method or winding number was run; the pencil independently returns the same $53$ and no further pole on the amplitude profile, which supports the count without proving it), and the real parts (the $4.7\times10^{-9}$ is an estimate; it says that no component of the window is off the line by more than the pencil's error, and it is not a proof of the Riemann hypothesis in the window: a zero outside the window, or hidden in a component of the profile, is not excluded). Nothing is certified off the line. The data enter the pencil as midpoints in double precision, an error of about $10^{-16}$, and the accuracy falls at the edge of the taper: the further $14$ poles kept at $6$ to $7.5$ from the centre, outside the window and not compared with the scan, have $\lvert\sigma-\tfrac12\rvert$ up to $2.6\times10^{-8}$. No priority is claimed: unrigorous computations of zeros at and far above this height exist \cite{Gourdon04,Odlyzko92,Odlyzko01}, and the sign change scan is the standard method; what this subsection records is that the deconvolution route of this chapter, from data on $\Re s=0.6$, reproduces the certified zeros of a window of this height with real parts read within $5\times10^{-9}$ of $\tfrac12$.\end{resultpar}

\section{Open problems}\label{es:sec:open}

\begin{problem}\label{es:prob:rigorous_pencil}
Rigorous pencil. The certification of Section~\ref{es:sec:certpencil} goes through the argument principle after the fact, and a verified pencil in the sense of Rump would need, besides verified linear algebra on ball data, a perturbation bound for the pencil under the model error of Theorem~\ref{es:thm:taper}, the far zeros and the remainder of Theorem~\ref{es:thm:taper}, which make the Hankel matrix of exact rank only in the limit. State and prove that bound, and compare the width of the resulting enclosures with the winding number route. Status: open.
\end{problem}

\begin{problem}\label{es:prob:spurious}
Spurious components. At the edges of overlapping windows a few components pass the amplitude filter without being zeros. Characterise them: they should be the pencil's representation of the taper's tail of the next zeros outside the window, with poles close to the genuine ones but amplitudes off the profile. A rule based on the agreement of two windows, or on the phase of the amplitude, which Theorem~\ref{es:thm:taper} predicts as $e^{i(\omega_0-a/\ell^2)(\gamma-c)}$, should remove all of them; under explicit hypotheses on the separation and the amplitudes, prove that it does. Status: open.
\end{problem}

\begin{problem}\label{es:prob:higher_wider}
Higher and wider. The runs near and above the verified range covered one window each, of fifty two and fifty three zeros, with data costing half an hour on two cores at $3\times10^{12}$ and $2.6$ seconds a point at $10^{13}$. The cost per point is the Riemann and Siegel sum, $t^{1/2}$ terms, and Hiary's method \cite{Hiary11} would bring it to $t^{1/3+o(1)}$; the number of points per zero is fixed by the taper and the spacing. Make the argument principle rigorous at large height, with an explicit bound for $|\zeta'|$ on a small disc from the explicit convexity bounds, so that the count in the window is certified and not only the existence on the line. Estimate the height at which a window of the pencil, with its full certificate, is feasible on a laptop in a day, and run it; and determine whether the many window scheme at a fixed height, say $10^{3}$ zeros above $3\,000\,175\,332\,800$, is competitive with Turing's method for that range when the real parts are wanted as output and not only tested. Status: open.
\end{problem}

\begin{problem}\label{es:prob:adaptive_prior}
The adaptive prior in full. Place the sine kernel process on marked distinct ordinates, a pair counting as one point marked by its displacement, and implement the posterior on the number, the positions and the marks of the zeros in a window, with the spike and slab on the marks, by reversible jump or by a determinantal birth and death sampler started at the pencil; measure on the many window run whether it removes the spurious components. State the count with its multiplicity (a double zero counts twice), and report the posterior probability of each count, conditional on the model, separately from the argument principle certificate, which it does not replace. Status: open.
\end{problem}

\begin{problem}\label{es:prob:constant_rate}
The constant in the rate. Theorem~\ref{es:thm:singular} gives $\delta^*=(z_\alpha\tau/\|t_\perp\|)^{1/2}$, with $\tau$ the standard deviation per real coordinate, and $\|t_\perp\|$ depends on the taper, the bandwidth and the level. Find the design that maximises $\|t_\perp\|/\tau$, the information and not $\|t_\perp\|$ alone, for a given cost, and determine how it scales with the height, through the gaps to the neighbouring zeros, which enter $t_\perp$ through the projection, and not through the mean density alone. Status: open.
\end{problem}
